\section{Two-block reduction and normalizations}\label{inf:sec-reduction}

Fix $p\ge64$ with $2p\in\Z$ and an integer $a$ with
$p/2\le a\le3p/4$. Write
\begin{equation}\label{inf:block-parameters}
 n=2p,\qquad b=2p-a,\qquad
 \ab=\frac a2=pr,\qquad \bb=\frac b2=p(1-r),
 \qquad \frac14\le r\le\frac38.
\end{equation}
For the spectral deformation argument the parameters $\ab,\bb$ will also
be allowed to be positive real numbers with sum $p$. All constructions
that define Euclidean domains are made at integer block sizes.
In particular $a,b\ge3$ and $\ab,\bb>1$ in this range.
We use $y=(X,Y)\in\R^a\times\R^b$ and, away from the origin, set
\begin{equation}\label{inf:coordinates}
 \varrho=|y|,\qquad t=|y|^2,\qquad
 x=\frac{|X|^2}{|X|^2+|Y|^2}.
\end{equation}
Here $t$ always denotes squared radius. A nonsquared reference radial
coordinate used later will be denoted by $\zeta$.

\begin{lemma}[Invariant differential operator]\label{inf:reduction}
For a twice differentiable invariant function $u(\varrho,x)$,
\begin{equation}\label{inf:invariant-laplacian}
 \Delta u=u_{\varrho\varrho}+\frac{2p-1}{\varrho}u_\varrho
 +\frac4{\varrho^2}\left\{x(1-x)u_{xx}+(\ab-px)u_x\right\}.
\end{equation}
Normalized spherical measure induces on $x$ the probability measure
\begin{equation}\label{inf:beta-measure}
 d\mu_r(x)=\frac{x^{\ab-1}(1-x)^{\bb-1}}{B(\ab,\bb)}\,dx.
\end{equation}
Normalized invariant volume measure on the unit ball is
$p t^{p-1}\,dt\,d\mu_r(x)$. Each angular polynomial occurs with
multiplicity one in this invariant sector.
\end{lemma}
\begin{proof}
Put $u(y)=v(q_X,q_Y)$, where $q_X=|X|^2$ and $q_Y=|Y|^2$.
Differentiating in each Euclidean coordinate and summing gives
\[
 \Delta u=4q_Xv_{q_Xq_X}+2av_{q_X}
          +4q_Yv_{q_Yq_Y}+2bv_{q_Y}.
\]
Use $q_X=\varrho^2x$ and $q_Y=\varrho^2(1-x)$ in this identity.
The mixed second derivatives in $\varrho,x$ cancel, and the remaining
coefficients are those in \eqref{inf:invariant-laplacian}.
Alternatively, the coefficients of $u_{xx}$ and $u_x$ follow directly
from $|\nabla_S x|^2=4x(1-x)$ and
$\Delta_S x=4(\ab-px)$.

Polar integration in the two blocks has volume element proportional to
$q_X^{\ab-1}q_Y^{\bb-1}\,dq_X\,dq_Y$. The determinant of the change
$(t,x)\mapsto(tx,t(1-x))$ has absolute value $t$. Its angular integral
is $B(\ab,\bb)$; its radial integral over $0<t<1$ is $1/p$.
Division by these two integrals proves the measure assertions.
An invariant function is a function of $\varrho,x$, so each angular
polynomial below occurs once.
\end{proof}

\subsection{The angular polynomials}

Let $\mathsf H_j(x)$ be the monic degree-$j$ orthogonal polynomial for $\mu_r$,
and define
\begin{equation}\label{inf:angular-normalization}
 G_j(x)=\frac{\mathsf H_j(x)}{\mathsf H_j(1)},\qquad
 \eta_j^2=\int_0^1G_j(x)^2\,d\mu_r(x),\qquad
 P_j(x)=\eta_j^{-1}G_j(x).
\end{equation}
The sign of $P_j$ is chosen to give a positive leading coefficient.
The two uses of Jacobi polynomials in this paper are distinguished by
letters: $G_j$ is angular, whereas $J_s^{\beta}$ below is radial.
The symbol $J_\nu$ without an upper index denotes the Bessel function.

\begin{lemma}[Angular formulas]\label{inf:angular-formulas}
For $j\ge0$,
\begin{align}
 G_j(x)&=\frac{P_j^{(\bb-1,\ab-1)}(2x-1)}{(\bb)_j/j!},
 &\mathsf H_j(1)&=\frac{(\bb)_j}{(p+j-1)_j},\label{inf:angular-jacobi}\\
 \eta_j^2&=\frac{(p-1)(\ab)_j j!}
 {(2j+p-1)(p-1)_j(\bb)_j},
 &\Lambda_rG_j&=\lambda_jG_j,\label{inf:angular-norm}\\
 \Lambda_r&=-4x(1-x)\partial_x^2-4(\ab-px)\partial_x,
 &\lambda_j&=4j(j+p-1).\nonumber
\end{align}
The convention for a rising factorial of length zero is $(z)_0=1$.
In the orthonormal basis multiplication by $x$ is
\begin{equation}\label{inf:jacobi-matrix}
 xP_j=\mathfrak a_{j+1}P_{j+1}+\mathfrak b_jP_j
       +\mathfrak a_jP_{j-1},
\end{equation}
where $\mathfrak a_0=0$ and
\begin{align}
 \mathfrak a_j^2&=
 \frac{j(j+\ab-1)(j+\bb-1)(j+p-2)}
 {(2j+p-2)^2(2j+p-1)(2j+p-3)},\quad j\ge1,
 \label{inf:recurrence-a}\\
 \mathfrak b_j&=\frac12+
 \frac{(r-1/2)p(p-2)}{(2j+p-2)(2j+p)}.
 \label{inf:recurrence-b}
\end{align}
All these identities hold also for real positive $\ab,\bb$ with $p\ge64$.
\end{lemma}
\begin{proof}
The shifted Rodrigues identity is
\[
 x^{\ab-1}(1-x)^{\bb-1}P_j^{(\bb-1,\ab-1)}(2x-1)
 =\frac{(-1)^j}{j!}\frac{d^j}{dx^j}
       \{x^{\ab+j-1}(1-x)^{\bb+j-1}\}.
\]
Repeated integration by parts against a polynomial of degree less than
$j$ proves orthogonality. The endpoint value is $(\bb)_j/j!$ and the
leading coefficient is $(p+j-1)_j/j!$. Their quotient proves the
second identity in \eqref{inf:angular-jacobi}. The integrations by parts
are valid: at each intermediate stage the endpoint powers in the
boundary terms have positive exponent. Differentiation of the Rodrigues
identity gives the displayed angular differential equation.

A further integration by parts against the degree-$j$ polynomial gives
\[
 \int_0^1\bigl(P_j^{(\bb-1,\ab-1)}(2x-1)\bigr)^2\,d\mu_r
 =\frac{(p-1)(\ab)_j(\bb)_j}
        {(2j+p-1)(p-1)_j j!}.
\]
Dividing by the squared endpoint value proves
\eqref{inf:angular-norm}. Orthogonality makes multiplication by $x$
tridiagonal and symmetric. The ratio of the leading coefficients of
$P_{j-1}$ and $P_j$, followed by the norm formula, gives
\eqref{inf:recurrence-a}. Comparing the coefficients of $x^j$ in the
monic recurrence gives \eqref{inf:recurrence-b}. These computations
hold for these positive real parameters.
\end{proof}

\begin{lemma}[Endpoint bound and angular derivatives]\label{inf:angular-sup}
For the parameters in \eqref{inf:block-parameters},
\[
 |G_j(x)|\le1\quad(0\le x\le1),\qquad 0<\eta_j\le1.
\]
On the unit sphere the homogeneous extension
$S_j(y)=|y|^{2j}G_j(x)$ satisfies
\begin{equation}\label{inf:angular-derivatives}
 \nrm{G_j}_{L^\infty(S^{2p-1})}\le1,\qquad
 \nrm{\nabla_S G_j}_\infty\le2j,\qquad
 \nrm{\nabla_S^2G_j}_{\mathrm{op},\infty}\le4j^2.
\end{equation}
The constants are independent of $p$ and of the admissible split.
\end{lemma}
\begin{proof}
For a Jacobi polynomial $v(z)=P_j^{(\alpha,\beta)}(z)$ with
$\alpha=\bb-1\ge\beta=\ab-1\ge0$ and $j>0$, write
$\kappa=j(j+\alpha+\beta+1)$. Its equation implies
\[
 \frac d{dz}\left(v^2+\frac{1-z^2}{\kappa}(v')^2\right)
 =\frac2\kappa\{\alpha-\beta+(\alpha+\beta+1)z\}(v')^2.
\]
The factor in braces is increasing. The expression in parentheses
therefore first decreases and then increases, with either interval
possibly empty. Its maximum is attained at an endpoint. The absolute
endpoint values are $(\alpha+1)_j/j!$ and $(\beta+1)_j/j!$; the former
is at least the latter. Division by the former proves the supremum
bound. The case $j=0$ is constant. The bound for $\eta_j$ follows by
integration against a probability measure.

The restriction to any great circle of the polynomial $S_j$ is a
trigonometric polynomial of degree at most $2j$. The first and second
Bernstein inequalities give bounds $2j$ and $4j^2$ for its derivatives
along that circle \cite[Vol.~II, Ch.~X, Theorem~3.13]{inf:zygmund}.
Great circles with prescribed unit tangent vector
compute the spherical gradient and the diagonal values of the
spherical Hessian. A real symmetric bilinear form has operator norm
equal to the supremum of its absolute diagonal values on unit vectors.
This proves \eqref{inf:angular-derivatives}.
\end{proof}

\begin{lemma}[Solid harmonics and radial profiles]\label{inf:solid-harmonics}
The function
\begin{equation}\label{inf:solid-definition}
 S_j(y)=t^jG_j(x)
\end{equation}
is an invariant homogeneous harmonic polynomial of degree $2j$.
For a twice differentiable radial profile $v$,
\begin{equation}\label{inf:profile-laplacian}
 \Delta\{S_jv(t)\}=4S_j\{tv''+(2j+p)v'\}.
\end{equation}
The regular solutions of $(\Delta+1)u=0$ in angular degree $2j$ are
multiples of
\begin{equation}\label{inf:regular-bessel-mode}
 \varrho^{1-p}J_{\nu_j}(\varrho)G_j(x),\qquad
 \nu_j=p-1+2j.
\end{equation}
Their apparent singularities at the origin are removable.
\end{lemma}
\begin{proof}
As $G_j$ is a polynomial of degree $j$, multiplication by $t^j$
expresses $S_j$ as a polynomial in $q_X,q_Y$. Applying
\eqref{inf:invariant-laplacian} and using its angular eigenvalue proves
harmonicity. For the product with $v(t)$, use
$\nabla t=2y$, $\Delta t=4p$, and $y\cdot\nabla S_j=2jS_j$.
The product rule gives \eqref{inf:profile-laplacian}.
Substituting $\varrho^{1-p}w(\varrho)$ in the radial Helmholtz equation
gives Bessel's equation of order $\nu_j$. Its regular solution is
$J_{\nu_j}$. Its convergent expansion is a constant times
$\varrho^{2j}$ multiplied by a power series in $\varrho^2$.
The factor $\varrho^{2j}G_j$ is the polynomial just constructed, which
proves regularity at the origin and across both block axes.
\end{proof}

\subsection{The ball field and the invariant Dirichlet spectrum}

Write $j_{\nu,k}$ for the $k$th positive zero of $J_\nu$ and set
$z_{j,k}=j_{\nu_j,k}$. All radial zero indices $k$ start at $1$.
Choose a positive zero $R=j_{p,k_0}$, where $k_0$ will be specified by
the arithmetic construction. The radial field with boundary value one is
\begin{equation}\label{inf:ball-field}
 u_0(\varrho)=
 \frac{\varrho^{1-p}J_{p-1}(\varrho)}{R^{1-p}J_{p-1}(R)}.
\end{equation}
Consecutive orders have no common positive zero: their recurrences
would give both the value and derivative of one solution equal to zero,
contradicting uniqueness for the Bessel equation. Thus the denominator
in \eqref{inf:ball-field} is nonzero.
The derivative identity
$(\varrho^{1-p}J_{p-1}(\varrho))'
=-\varrho^{1-p}J_p(\varrho)$ proves
\begin{equation}\label{inf:ball-traces}
 u_0(R)=1,\qquad u_0'(R)=0,\qquad u_0''(R)=-1.
\end{equation}
The last equality follows from the radial equation.

\begin{proposition}[Invariant Dirichlet realization]\label{inf:ball-spectrum}
The Friedrichs Dirichlet Laplacian $H_0=-\Delta_D$ in the invariant
subspace of the unit ball has eigenvalues $z_{j,k}^2$, one radial
sequence for each $j\ge0$, with repetitions retained. It has compact
resolvent. After the unitary radial substitution
$v(\zeta)=\sqrt{2p}\,\zeta^{p-1/2}u(\zeta)$ and the angular basis $P_j$, its
$j$th radial operator is
\begin{equation}\label{inf:radial-schrodinger}
 -\frac{d^2}{d\zeta^2}+\frac{\nu_j^2-1/4}{\zeta^2},
 \qquad 0<\zeta<1,
\end{equation}
with Friedrichs condition at zero and Dirichlet condition at one.
The spectrum does not depend on the split $r$.
\end{proposition}
\begin{proof}
The orthogonal group action commutes with the Dirichlet form. Averaging
over it is an orthogonal projection preserving $H^1_0$, so restriction
to the invariant subspace gives a closed Dirichlet form. Compactness of
$H^1_0(B)\hookrightarrow L^2(B)$ implies compact resolvent for this
restriction. The polynomials $P_j$ form a complete angular orthonormal
basis; this follows from density of polynomials in the continuous
functions on $[0,1]$ and density of continuous functions in $L^2(\mu_r)$.
Separation of variables gives \eqref{inf:radial-schrodinger}. Its
regular solution at energy $z^2$ is a multiple of
$\sqrt\zeta J_{\nu_j}(z\zeta)$. The other independent solution is not
in the Friedrichs form domain, as $\nu_j\ge p-1>1$.
The boundary condition at one is $J_{\nu_j}(z)=0$.
Neither $\nu_j$ nor the separated eigenvalue equation contains $r$.
\end{proof}

\begin{lemma}[Radial polynomial orthogonality]\label{inf:radial-polynomials}
Set $\beta_j=2j+p-1$ and
\begin{equation}\label{inf:radial-basis}
 J_s^\beta(t)=P_s^{(0,\beta)}(2t-1),\qquad
 \Psi_{j,s}(y)=S_j(y)J_s^{\beta_j}(t),\quad s\ge0.
\end{equation}
Then $J_s^\beta(1)=1$ and
\begin{align}
 J_s^\beta(t)&=\sum_{h=0}^s(-1)^{s-h}\binom sh
                         \binom{s+\beta+h}{s}t^h,
                         \label{inf:radial-expansion}\\
 \int_0^1t^\beta J_s^\beta J_v^\beta\,dt
 &=\frac{\delta_{sv}}{2s+\beta+1},\qquad
 \nrm{\Psi_{j,s}}_{L^2(B)}^2
 =\eta_j^2\frac p{p+2j+2s}.
 \label{inf:radial-norm}
\end{align}
Here volume is normalized as in Lemma~\ref{inf:reduction}. All invariant
polynomials are finite linear combinations of these basis elements.
\end{lemma}
\begin{proof}
Apply the Rodrigues computation of Lemma~\ref{inf:angular-formulas}
with Jacobi parameters $0,\beta$. It gives the coefficient formula,
endpoint value, and integral norm. Multiplication by $S_j=t^jG_j$
changes the radial weight $p t^{p-1}$ to $p t^{\beta_j}$; angular
orthogonality then proves the last identity.
A monomial in $q_X,q_Y$ is $t^d$ times a polynomial in $x$ of degree
at most $d$. Expand that polynomial in the $G_j$, $j\le d$, and expand
the remaining power $t^{d-j}$ in the $J_s^{\beta_j}$. This proves the
spanning assertion.
\end{proof}
